% !TEX program = pdflatex
%=====================================================================
%  Criptografía y Seguridad - ITBA
%  Clase 6 — Políticas de seguridad y Control de acceso
%  Apunte integral: teoría + práctica
%=====================================================================
\documentclass[11pt,a4paper]{article}

%---------------------------------------------------------------------
%  Paquetes
%---------------------------------------------------------------------
\usepackage[utf8]{inputenc}
\usepackage[T1]{fontenc}
\usepackage[spanish,es-tabla,es-noshorthands]{babel}
\usepackage{lmodern}
\usepackage{microtype}

\usepackage{amsmath,amssymb,amsthm}
\usepackage{mathtools}
\usepackage{geometry}
\geometry{a4paper,margin=2.4cm,headheight=22pt}

\usepackage{xcolor}
\usepackage{graphicx}
\usepackage{enumitem}
\usepackage{booktabs}
\usepackage{array}
\usepackage{tcolorbox}
\tcbuselibrary{skins,breakable}
\usepackage{fancyhdr}
\usepackage{titlesec}
\usepackage{hyperref}
\usepackage{listings}
\usepackage{caption}

\usepackage{tikz}
\usetikzlibrary{shapes.geometric,shapes.misc,shapes.symbols,arrows.meta,positioning,calc,
                fit,backgrounds,decorations.pathmorphing,decorations.pathreplacing,
                shadows.blur,chains,matrix}

%---------------------------------------------------------------------
%  Paleta de color (inspirada en las diapositivas de la cátedra)
%---------------------------------------------------------------------
\definecolor{itbaCyan}{HTML}{6EC6D9}
\definecolor{itbaCyanDark}{HTML}{2E8B9E}
\definecolor{itbaBlue}{HTML}{AEDCEA}
\definecolor{boxBlue}{HTML}{BBD4F0}
\definecolor{slideGray}{HTML}{ECECEC}
\definecolor{practGreen}{HTML}{4F8A3D}
\definecolor{practGreenL}{HTML}{E2EFD9}
\definecolor{attackRed}{HTML}{C0392B}
\definecolor{codeBg}{HTML}{F5F5F5}
% niveles de clasificación (estilo slides Bell-LaPadula / Biba)
\definecolor{lvlTS}{HTML}{E74C3C}
\definecolor{lvlS}{HTML}{F39C12}
\definecolor{lvlC}{HTML}{5DADE2}
\definecolor{lvlU}{HTML}{58D68D}
% MAC verde / DAC amarillo (resaltados de la práctica)
\definecolor{macGreen}{HTML}{2ECC71}
\definecolor{dacYellow}{HTML}{F7DC6F}

%---------------------------------------------------------------------
%  Hipervínculos
%---------------------------------------------------------------------
\hypersetup{
  colorlinks=true,
  linkcolor=itbaCyanDark,
  urlcolor=itbaCyanDark,
  citecolor=itbaCyanDark,
  pdftitle={Clase 6 - Politicas y Control de acceso},
  pdfauthor={Criptografia y Seguridad - ITBA}
}

%---------------------------------------------------------------------
%  Encabezados y pies
%---------------------------------------------------------------------
\pagestyle{fancy}
\fancyhf{}
\lhead{\small\textsc{Criptografía y Seguridad — ITBA}}
\rhead{\small Clase 6 · Políticas y Control de acceso}
\cfoot{\small\thepage}
\renewcommand{\headrulewidth}{0.4pt}
\renewcommand{\headrule}{\hbox to\headwidth{\color{itbaCyanDark}\leaders\hrule height \headrulewidth\hfill}}

%---------------------------------------------------------------------
%  Títulos de sección con barra de color (estilo diapositiva)
%---------------------------------------------------------------------
\titleformat{\section}
  {\normalfont\Large\bfseries\color{itbaCyanDark}}
  {\thesection}{0.6em}{}
  [\vspace{-0.4em}{\color{itbaCyan}\titlerule[2pt]}]
\titleformat{\subsection}
  {\normalfont\large\bfseries\color{itbaCyanDark}}
  {\thesubsection}{0.6em}{}
\titleformat{\subsubsection}
  {\normalfont\normalsize\bfseries\color{black}}
  {\thesubsubsection}{0.6em}{}

%---------------------------------------------------------------------
%  Cajas reutilizables
%---------------------------------------------------------------------
% Caja "diapositiva" (recuadro gris de las slides)
\newtcolorbox{slidebox}[1][]{
  enhanced, colback=slideGray, colframe=black!55, boxrule=0.5pt,
  arc=1pt, left=6pt,right=6pt,top=4pt,bottom=4pt, #1}

% Caja de definición
\newtcolorbox{defbox}[1]{
  enhanced, breakable, colback=itbaBlue!25, colframe=itbaCyanDark,
  boxrule=1pt, arc=2pt, left=8pt,right=8pt,top=6pt,bottom=6pt,
  fonttitle=\bfseries, coltitle=white, title={#1}}

% Caja de nota práctica
\newtcolorbox{practbox}[1]{
  enhanced, breakable, colback=practGreenL, colframe=practGreen,
  boxrule=1pt, arc=2pt, left=8pt,right=8pt,top=6pt,bottom=6pt,
  fonttitle=\bfseries, coltitle=white, title={#1}}

% Caja de atención / ataque
\newtcolorbox{warnbox}[1]{
  enhanced, breakable, colback=attackRed!8, colframe=attackRed,
  boxrule=1pt, arc=2pt, left=8pt,right=8pt,top=6pt,bottom=6pt,
  fonttitle=\bfseries, coltitle=white, title={#1}}

% Caja de anécdota / aclaración del profe
\newtcolorbox{anecbox}[1]{
  enhanced, breakable, colback=boxBlue!25, colframe=itbaCyan!80!black,
  boxrule=0.8pt, arc=2pt, left=8pt,right=8pt,top=6pt,bottom=6pt,
  fonttitle=\bfseries\itshape, coltitle=black, title={#1}}

%---------------------------------------------------------------------
%  Estilo de código
%---------------------------------------------------------------------
\lstdefinestyle{code}{
  basicstyle=\ttfamily\footnotesize,
  backgroundcolor=\color{codeBg},
  keywordstyle=\color{itbaCyanDark}\bfseries,
  commentstyle=\color{practGreen}\itshape,
  stringstyle=\color{attackRed},
  breaklines=true, frame=single, rulecolor=\color{black!30},
  showstringspaces=false, columns=fullflexible, tabsize=2,
  xleftmargin=10pt, framexleftmargin=8pt}

%---------------------------------------------------------------------
%  Macros matemáticos
%---------------------------------------------------------------------
\newcommand{\Sset}{\mathcal{S}}
\newcommand{\Oset}{\mathcal{O}}
\newcommand{\Rset}{\mathcal{R}}
\newcommand{\dom}{\mathrel{\mathbf{dom}}}        % relación de dominancia
\newcommand{\il}{\mathrm{il}}                     % integrity level
\newcommand{\Lvl}{L}                              % security level
\newcommand{\highlight}[2]{\colorbox{#1}{\strut #2}}

%=====================================================================
\begin{document}
%=====================================================================

%---------------------------------------------------------------------
%  Portada
%---------------------------------------------------------------------
\begin{titlepage}
\centering
\vspace*{1.2cm}

\begin{tikzpicture}
  \node[rounded corners=2pt, fill=itbaCyan, text=black, inner sep=14pt,
        minimum width=0.8\textwidth, align=center] (t)
        {\Huge\bfseries Criptografía y Seguridad};
\end{tikzpicture}

\vspace{0.6cm}
{\LARGE Seguridad en sistemas y aplicaciones}\\[0.25cm]
{\Huge\bfseries\color{itbaCyanDark} Políticas de seguridad\\[2pt] y Control de acceso}

\vspace{0.8cm}

% Ilustración: bóveda / caja fuerte (recreación del icono de portada)
\begin{tikzpicture}[scale=1.0]
  \fill[black!12] (-3.2,-2.4) rectangle (3.2,2.8);
  \draw[line width=2pt, black!55] (-3.2,-2.4) rectangle (3.2,2.8);
  \fill[itbaCyanDark!75] (-2.4,-2.0) rectangle (2.4,2.4);
  \draw[line width=3pt, black!70] (-2.4,-2.0) rectangle (2.4,2.4);
  \draw[line width=1.2pt, itbaCyan] (-2.0,-1.6) rectangle (2.0,2.0);
  \draw[line width=2pt, black!70] (0,-2.0) -- (0,2.4);
  \foreach \a in {0,45,...,315}{
    \draw[line width=2pt, black!70] (-1.1,0.2) -- ++(\a:0.6);
  }
  \draw[line width=2.5pt, black!70] (-1.1,0.2) circle (0.42);
  \fill[itbaCyan] (-1.1,0.2) circle (0.16);
  \draw[line width=2.5pt, black!70] (1.1,0.2) circle (0.42);
  \fill[itbaCyan] (1.1,0.2) circle (0.16);
  \foreach \y in {-1.2,-0.4,0.8,1.6}{
     \fill[black!70] (-2.2,\y) circle (0.07);
     \fill[black!70] (2.2,\y) circle (0.07);
  }
\end{tikzpicture}

\vfill
{\large Apunte integral — Teoría + Práctica}\\[0.2cm]
{\normalsize Clase 6 \quad·\quad Capítulos 4--8, \emph{Computer Security: Art and Science} (M.~Bishop)}\\[0.4cm]
{\small Instituto Tecnológico de Buenos Aires (ITBA)}

\vspace{0.6cm}
{\footnotesize\itshape Documento elaborado a partir de la transcripción de la clase teórica
(Prof.\ Pablo Abad) y de las diapositivas oficiales de la teórica y de la práctica.}
\end{titlepage}

\tableofcontents
\newpage

%=====================================================================
\section{Definiendo seguridad en aplicaciones}
%=====================================================================

En criptografía aprendimos algo fundamental: \textbf{no hay una única definición de
seguridad}. Para una misma función criptográfica podía haber distintas pruebas, y cada prueba
era, en el fondo, \emph{una definición distinta de seguridad}. Al pasar de funciones
matemáticas a \textbf{sistemas dinámicos}, ocurre lo mismo: la seguridad depende del contexto
(aplicación, servicio, empresa).

\begin{defbox}{Política de seguridad — la definición operativa}
Existe una definición genérica para sistemas complejos: si pudiésemos clasificar todos los
\textbf{estados} de un sistema en \emph{seguros} e \emph{inseguros}, entonces
\begin{center}\itshape
un sistema es seguro mientras sus estados queden dentro de la política de seguridad,\\
es decir, mientras transicione únicamente entre estados que consideramos seguros.
\end{center}
La \textbf{política de seguridad} es, literalmente, \emph{la definición de qué estados de un
sistema vamos a considerar seguros} frente al resto.
\end{defbox}

Esta definición es más teórica que práctica —si modelamos una aplicación como una máquina de
estados, la cantidad de estados es \emph{abismalmente grande} y difícil de clasificar—, pero
aporta dos conceptos muy poderosos que sí aplican siempre:

\begin{enumerate}[leftmargin=1.6em,itemsep=2pt]
  \item Necesitamos una forma de \textbf{caracterizar}, sacando ``una foto en el tiempo'',
        si el estado actual del sistema es seguro o no. \emph{Eso es la política de seguridad.}
  \item Un sistema que pasa \textbf{eventualmente} a un estado inseguro \textbf{deja de ser
        seguro}, \emph{aunque después vuelva} a un estado que consideremos seguro.
\end{enumerate}

% ---- Diagrama: estados seguros / inseguros (recreación slide 2) ----
\begin{figure}[h]
\centering
\begin{tikzpicture}[font=\small,>=Stealth]
  % región "política de seguridad"
  \draw[thick,black!55,fill=itbaBlue!18] plot[smooth cycle] coordinates
        {(-3.4,0) (-2.8,1.5) (-1,1.9) (1.2,1.7) (1.9,0.4) (1.1,-1.4) (-1.2,-1.8) (-3,-1.1)};
  \node[font=\scriptsize\itshape,black!60] at (-2.2,1.55) {estados seguros};

  % nodos seguros (azules)
  \node[circle,draw=black,line width=1.4pt,fill=lvlC,minimum size=0.42cm] (a) at (-2.3,0.3) {};
  \node[circle,draw,fill=lvlC,minimum size=0.38cm] (b) at (-0.9,1.0) {};
  \node[circle,draw,fill=lvlC,minimum size=0.38cm] (c) at (-0.5,-0.2) {};
  \node[circle,draw,fill=lvlC,minimum size=0.38cm] (d) at (-1.6,-1.1) {};
  \node[circle,draw,fill=lvlC,minimum size=0.38cm] (e) at (1.0,0.5) {};
  % nodos inseguros (rojos)
  \node[circle,draw,fill=lvlTS,minimum size=0.38cm] (f) at (3.0,1.0) {};
  \node[circle,draw,fill=lvlTS,minimum size=0.38cm] (g) at (3.4,-0.6) {};
  \node[circle,draw,fill=lvlTS,minimum size=0.38cm] (h) at (4.5,-1.3) {};
  \node[circle,draw,fill=lvlTS,minimum size=0.38cm] (i) at (1.4,-1.6) {};
  \node[font=\scriptsize\itshape,attackRed] at (3.9,1.0) {estados inseguros};

  \draw[black!70] (a)--(b) (a)--(c) (a)--(d) (c)--(d) (c)--(e) (b)--(e);
  \draw[dashed,black!55] (e)--(f) (e)--(i) (d)--(i);
  \draw[black!70] (f)--(g) (g)--(h) (i)--(g);
\end{tikzpicture}
\caption{Mientras el sistema transita por estados \emph{dentro} de la región (azules), es
seguro. Si una transición lo lleva a un estado inseguro (rojo), deja de ser seguro.
\textit{(Recreación de la diapositiva 2.)}}
\end{figure}

\begin{anecbox}{El costo de volver a un estado seguro}
El punto~2 tiene enorme importancia en seguridad defensiva. Si a un servidor se le
\textbf{comprometió la seguridad}, las buenas prácticas dicen que habría que ``prenderle
fuego y traer uno nuevo'' (o, sin tanto drama físico, \textbf{reinstalarlo}). Es muy difícil,
una vez comprometida, volver a un estado del que se tenga garantía. Por eso trabajar en
seguridad requiere cierto \emph{mindset} de adaptabilidad y resiliencia: como en criptografía,
no se puede bajar la probabilidad de incidente a cero, y cuando hay un incidente el costo de
recuperación es altísimo.
\end{anecbox}

\subsection{La tríada CIA}
La seguridad puede enfocarse desde varios lados. En criptografía vimos
\textbf{confidencialidad} y \textbf{integridad}; el espectro completo incorpora un tercer
servicio, la \textbf{disponibilidad}. A estos tres se los conoce como la \textbf{tríada CIA}
(\emph{Confidentiality, Integrity, Availability}).

\begin{defbox}{Tríada CIA}
\begin{itemize}[leftmargin=1.5em,itemsep=2pt]
  \item \textbf{Confidencialidad}: atañe a la \emph{diseminación} de la información. Hay
        información que debe ser accesible a cierto grupo y no a otro. (\emph{Guardar la
        información en secreto: acceso solo a los autorizados.})
  \item \textbf{Integridad}: tiene que ver con relaciones de \emph{confianza}. Necesito poder
        certificar la calidad de algo: que una transacción ocurrió, que los datos son
        correctos, o que el \emph{origen} es quien dice ser. Se distingue la
        \textbf{integridad de datos} de la \textbf{integridad de origen}.
  \item \textbf{Disponibilidad}: capacidad de acceder a la información o ejecutar una acción
        \emph{en el momento} en que se la necesita. Involucra \textbf{temporalidad}, es más
        dinámica y por eso no se estudia tanto en criptografía (que es estática).
\end{itemize}
\end{defbox}

% ---- Diagrama de la tríada (recreación slide 3 - diagrama de Venn) ----
\begin{figure}[h]
\centering
\begin{tikzpicture}[font=\small]
  \begin{scope}[blend group=multiply]
    \fill[itbaCyanDark!70]  (0,0.9)   circle (1.6);
    \fill[itbaCyanDark!45]  (-1.0,-0.7) circle (1.6);
    \fill[itbaCyan!75]      (1.0,-0.7)  circle (1.6);
  \end{scope}
  \node[white,font=\bfseries,align=center] at (0,1.5) {Confiden-\\cialidad};
  \node[white,font=\bfseries] at (-1.6,-1.1) {Integridad};
  \node[white,font=\bfseries] at (1.6,-1.1) {Disponib.};
  \node[font=\scriptsize\bfseries,white] at (0,-0.15) {C·I·A};
\end{tikzpicture}
\caption{La tríada CIA. La seguridad de un activo se preserva sosteniendo las tres
propiedades. \textit{(Recreación de la diapositiva 3.)}}
\end{figure}

%=====================================================================
\section{Conceptos de control de acceso}
%=====================================================================

Para construir teoría sobre seguridad se introducen tres conceptos básicos, engañosamente
simples, que son la base de todo:

\begin{defbox}{Sujeto, Objeto y Acceso}
\begin{description}[leftmargin=2.0em,style=nextline,itemsep=2pt]
  \item[Sujeto] La entidad que es \emph{origen} de algún estímulo: ejecuta acciones sobre
        otras entidades. Pueden ser personas, servicios, aplicaciones, servidores.
  \item[Objeto] El \emph{destino} de esos estímulos. Generalmente recursos: datos, servicios,
        bases de datos, tablas, programas, archivos.
  \item[Acceso] Las \emph{formas de interacción} entre un sujeto y un objeto (las operaciones
        posibles: leer, escribir, borrar, mover, ejecutar, \dots).
\end{description}
\end{defbox}

% ---- Diagrama sujeto-objeto-acceso (recreación slide 5) ----
\begin{figure}[h]
\centering
\begin{tikzpicture}[font=\small,>=Stealth]
  \node[circle,draw,fill=itbaBlue,minimum size=0.9cm] (s) at (0,0) {};
  \node[font=\scriptsize] at (0,-0.85) {\textbf{Sujeto}};
  \node[draw,fill=black!12,minimum width=0.9cm,minimum height=0.9cm] (o) at (4,0) {};
  \node[font=\scriptsize] at (4,-0.85) {\textbf{Objeto}};
  \draw[->,line width=1.2pt] (s) -- (o);
  \draw[->,attackRed,line width=1.2pt,bend left=40] (2.7,0.9) to (2.0,0.18);
  \node[attackRed,font=\scriptsize\bfseries] at (2.0,1.15) {Acceso};
\end{tikzpicture}
\caption{Un sujeto interactúa con un objeto mediante un acceso. \textit{(Recreación de la
diapositiva 5.)}}
\end{figure}

\begin{practbox}{El ejemplo de siempre: Linux}
En Linux (paradigma Unix) el esquema es directo: los \textbf{usuarios} son los sujetos, los
\textbf{archivos} son los objetos (``todo es un archivo'') y los \textbf{permisos} son los
accesos que rigen qué operaciones se pueden hacer. Es un modelo \emph{engañosamente simple}:
todos interactuamos con uno así, pero esconde mucha complejidad (lo retomamos en la
sección de ejemplos).
\end{practbox}

\subsection{La tríada desde el punto de vista de los accesos}
La importancia de la tríada \emph{sujeto–objeto–acceso} es que permite reducir los tres
servicios de seguridad a una mínima expresión modelable:

\begin{slidebox}
\centering\itshape Las políticas de seguridad se implementan controlando cómo los sujetos
acceden a los objetos.
\end{slidebox}

\begin{itemize}[leftmargin=1.5em,itemsep=2pt]
  \item \textbf{Confidencialidad}: existen \underline{sujetos} que \underline{NO pueden
        acceder} a ciertos \underline{objetos}.
  \item \textbf{Integridad}: existen \underline{objetos} que \underline{solo pueden ser
        modificados} por ciertos \underline{sujetos confiables}.
  \item \textbf{Disponibilidad}: existen \underline{objetos} que pueden ser
        \underline{accedidos} por ciertos \underline{sujetos} cuando lo requieran.
\end{itemize}

\subsection{Identidad del sujeto}
\textbf{La identidad es lo que identifica unívocamente al sujeto} dentro del universo de
posibles sujetos. Es un \emph{requerimiento para la autenticación} (tema de la clase
siguiente). Ejemplos: en Windows/Active Directory el \texttt{SID}
(\texttt{S-1-5-21-\dots}) y el \texttt{DistinguishedName}; en Linux el \texttt{UID} (la
primera columna numérica de \texttt{/etc/passwd}, p.\ ej.\ \texttt{root:x:0:0:\dots}).

%=====================================================================
\section{Políticas vs.\ Mecanismos}
%=====================================================================

\begin{practbox}{Distinción clave (clase práctica)}
\begin{center}
\renewcommand{\arraystretch}{1.3}
\begin{tabular}{>{\raggedright\arraybackslash}p{6.4cm} >{\raggedright\arraybackslash}p{6.4cm}}
\toprule
{\color{attackRed}\bfseries Políticas} & {\color{itbaCyanDark}\bfseries Mecanismos}\\
\midrule
Definen \textbf{qué acciones son seguras} (permitidas) y \textbf{qué acciones son inseguras}
(prohibidas). &
Procedimiento, herramienta o método para \textbf{garantizar el cumplimiento} de la política.\\
\bottomrule
\end{tabular}
\end{center}
\medskip
Los mecanismos cumplen tres funciones complementarias:
\begin{center}
\begin{tikzpicture}[font=\small,>=Stealth,node distance=6mm]
  \node[draw,rounded corners,fill=boxBlue!40] (p) {Mecanismo};
  \node[draw,rounded corners,fill=practGreenL,below left=8mm and 6mm of p] (prev) {Prevenir};
  \node[draw,rounded corners,fill=practGreenL,below=10mm of p] (det) {Detectar};
  \node[draw,rounded corners,fill=practGreenL,below right=8mm and 6mm of p] (rec) {Recuperar};
  \draw[->] (p)--(prev); \draw[->] (p)--(det); \draw[->] (p)--(rec);
\end{tikzpicture}
\end{center}
\end{practbox}

\noindent En síntesis: la \textbf{política} dice qué se considera seguro; el
\textbf{mecanismo} es el cómo se hace cumplir.

%=====================================================================
\section{Modelos de seguridad}
%=====================================================================

Tenemos la idea de sujeto–objeto y la de servicios (confidencialidad, etc.), pero ante un
problema práctico necesitamos resolver \emph{cómo} definir qué estados son seguros.

\begin{defbox}{¿Qué es un modelo de seguridad?}
Un \textbf{modelo de seguridad} es un \textbf{modelo formal y lógico} que define las
\textbf{reglas e interacciones} que permiten implementar un mecanismo de control de acceso.
A través de una sustentación formal (matemática) se demuestra que el sistema es seguro.
\begin{itemize}[leftmargin=1.4em,itemsep=1pt]
  \item Están \textbf{especializados}: hay modelos de \emph{confidencialidad}, de
        \emph{integridad}, e incluso de \emph{disponibilidad}.
  \item Son \textbf{generalizaciones} que ya sabemos que funcionan: evitan
        \emph{reinventar la rueda} y razonar a un nivel más alto (analogía: el modelo es a un
        lenguaje de alto nivel lo que la solución artesanal es al ensamblador).
  \item \textbf{No alcanzan por sí solos} en un sistema real: se aplican a \emph{partes} del
        sistema.
\end{itemize}
\end{defbox}

Hay muchísimos modelos, pero dos o tres dominan el escenario y conviene conocerlos:
\textbf{Bell--LaPadula} (confidencialidad), \textbf{Biba} (integridad) y
\textbf{Clark--Wilson} (integridad en entornos comerciales).

%=====================================================================
\section{Modelo Bell--LaPadula (confidencialidad)}
%=====================================================================

Bell y LaPadula (investigadores de IBM) diseñaron este modelo, célebre por ser el
\textbf{primer modelo de políticas de confidencialidad aprobado para uso en organismos
gubernamentales de EE.UU.} y el primero \emph{público} utilizable en ambientes civiles
(análogo, en su impacto, a lo que fue DES). Está pensado para el \textbf{ámbito militar},
focalizado en proteger la \textbf{confidencialidad}.

\begin{defbox}{Objetivo y elementos}
\textbf{Objetivo:} prevenir el \emph{acceso no autorizado} a la información; las
modificaciones no autorizadas son \emph{secundarias}. Combina \highlight{macGreen}{acceso
mandatorio} y \highlight{dacYellow}{acceso discrecional}. Elementos:
\begin{itemize}[leftmargin=1.4em,itemsep=1pt]
  \item Sujetos y Objetos.
  \item Modos de acceso $=\{\underline{r}\text{ead}, \underline{w}\text{rite},\dots\}$.
  \item Clasificación de seguridad $=\{\text{TS}, \text{S}, \text{C}, \text{U}\dots\}$
        (\emph{Top Secret, Secret, Confidential, Unclassified}).
  \item Niveles de seguridad $=$ Clasificación $\times$ Categoría (reticulado / \emph{lattice}).
\end{itemize}
La información \textbf{debe clasificarse} según su nivel de confidencialidad, y todos los
sujetos deben tener un \textbf{nivel de acceso} definido, con los mismos posibles niveles que
la información.
\end{defbox}

% ---- Diagrama: niveles (recreación slide 11) ----
\begin{figure}[h]
\centering
\begin{tikzpicture}[font=\small,node distance=3mm]
  \node[draw,rounded corners,fill=lvlTS,minimum width=3.2cm,minimum height=0.8cm,text=white,font=\bfseries] (ts) {Top Secret};
  \node[draw,rounded corners,fill=lvlS,minimum width=3.2cm,minimum height=0.8cm,below=of ts] (s) {Secret};
  \node[draw,rounded corners,fill=lvlC,minimum width=3.2cm,minimum height=0.8cm,below=of s] (c) {Confidential};
  \node[draw,rounded corners,fill=lvlU,minimum width=3.2cm,minimum height=0.8cm,below=of c] (u) {Unclassified};
  \draw[<->,line width=1.2pt,itbaCyanDark] (4.4,-2.45) -- node[right,font=\scriptsize,align=left]{mayor\\ nivel} (4.4,0.45);
\end{tikzpicture}
\caption{Niveles de clasificación del modelo original. \textit{(Recreación de la
diapositiva 11.)}}
\end{figure}

\subsection{Reglas (principios)}
El modelo tiene dos reglas, más una tercera que sintetiza a las dos. Para enunciarlas
con generalidad usamos la \textbf{relación de dominancia} $\dom$ y notamos $\Lvl(s)$ y
$\Lvl(o)$ los niveles del sujeto $s$ y del objeto $o$.

\begin{defbox}{Las tres reglas de Bell--LaPadula}
\begin{enumerate}[leftmargin=1.7em,itemsep=4pt]
  \item \textbf{Condición de seguridad simple — \emph{Read Down}.} El sujeto $s$ puede
        \emph{leer} el objeto $o$ si y solo si
        \[ \Lvl(s) \dom \Lvl(o) \quad\text{y } s \text{ tiene permiso para leer } o. \]
        Un sujeto puede leer objetos de su nivel \emph{hacia abajo}.
  \item \textbf{Propiedad $\ast$ (estrella) — \emph{Write Up}.} El sujeto $s$ puede
        \emph{escribir} el objeto $o$ si y solo si
        \[ \Lvl(o) \dom \Lvl(s) \quad\text{y } s \text{ tiene permiso para escribir } o. \]
        Un sujeto solo puede escribir hacia su nivel \emph{o más arriba}.
  \item \textbf{Propiedad estrella fuerte (síntesis).} Si un sujeto necesita
        \emph{leer y escribir} el mismo objeto, solo podrá hacerlo con objetos de su
        \textbf{misma clasificación} (es lo único que satisface ambas reglas a la vez).
\end{enumerate}
\textbf{Regla mnemotécnica:} \emph{Read down, Write up.} Pensar ``escritura'' y ``lectura''
como operaciones que \emph{cambian} o \emph{no cambian} el estado: leer puede ser ver
contenido, metadata, tamaño o existencia; escribir puede ser modificar, borrar, agregar.
\end{defbox}

% ---- Diagrama: read down / write up (recreación slides 12 / práctica) ----
\begin{figure}[h]
\centering
\begin{tikzpicture}[font=\small,>=Stealth]
  % flecha de flujo de información
  \draw[->,line width=6pt,black!15] (7.6,-2.2) -- node[right=3pt,black!55,font=\scriptsize]{Flujo Info} (7.6,2.2);
  % Read down
  \node[circle,draw,fill=itbaBlue,minimum size=0.7cm] (sr) at (0,1.2) {};
  \node[draw,fill=black!10,minimum size=0.55cm] (or1) at (3,2.0) {};
  \node[draw,fill=black!10,minimum size=0.55cm] (or2) at (3,0.4) {};
  \draw[<-,line width=1pt] (sr) -- node[above,font=\scriptsize]{r} (or1);
  \draw[<-,line width=1pt] (sr) -- node[below,font=\scriptsize]{r} (or2);
  \node[font=\scriptsize\bfseries,itbaCyanDark] at (1.5,2.6) {Read Down: $L(s)\dom L(o)$};
  % Write up
  \node[circle,draw,fill=itbaBlue,minimum size=0.7cm] (sw) at (0,-1.6) {};
  \node[draw,fill=black!10,minimum size=0.55cm] (ow1) at (3,-0.8) {};
  \node[draw,fill=black!10,minimum size=0.55cm] (ow2) at (3,-2.4) {};
  \draw[->,line width=1pt] (sw) -- node[above,font=\scriptsize]{w} (ow1);
  \draw[->,line width=1pt] (sw) -- node[below,font=\scriptsize]{w} (ow2);
  \node[font=\scriptsize\bfseries,itbaCyanDark] at (1.5,-3.0) {Write Up: $L(o)\dom L(s)$};
\end{tikzpicture}
\caption{Un sujeto lee hacia abajo y escribe hacia arriba; la información ``sube''.
\textit{(Recreación de las diapositivas de teórica y práctica.)}}
\end{figure}

\begin{anecbox}{¿Por qué la regla de escritura? El control de \emph{flujo} de información}
La primera regla resuelve el \textbf{acceso directo}: si clasifico un documento como
\emph{Top Secret} y solo yo y mi socio somos TS, nadie más debería leerlo. Pero supongamos
que la escritura fuese libre. Abro el documento en un editor que, por las dudas, guarda una
\emph{copia temporal} en otro directorio. Ese nuevo archivo tiene la misma información pero
\emph{otras reglas}: ¿quién garantiza que no lo lea alguien que no debe? La primera regla
sola no impide que alguien con acceso ``vuelque'' la información en otro lado. La
\textbf{regla de cierre (\emph{write up})} ataca eso: si tenés acceso a información muy
sensible, \emph{no podés escribir hacia niveles más bajos}, porque estarías filtrándola.
Bell--LaPadula no garantiza confidencialidad estática, sino un \textbf{control de flujo de
información}: hay demostraciones formales (sobre máquinas de Turing, en el libro de Bishop)
de que \emph{no existe ningún camino} por el que la información ``baje'' de nivel.
\end{anecbox}

\begin{anecbox}{El extremo: el Proyecto Manhattan}
El modelo es \textbf{muy estricto}: ``si accediste a este tipo de información, ya no podés
hablar con el resto de los mortales''. Esto importa en inteligencia militar, donde una fuga
puede costar vidas. Incluso \emph{sin querer}, quien conoce información de alto nivel puede
filtrarla parcialmente (una respuesta tangencial, un gesto). La aplicación de esta regla llevó
a que, durante el desarrollo de la bomba atómica, EE.UU.\ \textbf{aislara físicamente a los
científicos del Proyecto Manhattan en un pueblo entero} junto con sus familias, para que no
interactuaran con gente fuera del proyecto: estaban aplicando Bell--LaPadula.
\end{anecbox}

\begin{warnbox}{Limitación: el principio de tranquilidad}
El mayor problema de Bell--LaPadula es el \textbf{principio de tranquilidad}: asume que la
clasificación de sujetos y objetos \emph{no cambia con el tiempo}. Pero en la realidad sí
cambia: los empleados ascienden o son despedidos, y la información hoy confidencial puede ser
pública en unos meses. Las implementaciones reales agregan \textbf{mecanismos al costado para
reclasificar}. Un síntoma típico: los \emph{emails} de gente de organismos gubernamentales
llevan al pie una leyenda de \emph{desclasificación} (``\emph{declassified for use at
unclassified level}''), porque según las reglas estrictas un sujeto con acceso \emph{secret}
no debería poder interactuar en foros abiertos. El modelo \textbf{Clark--Wilson} se usa como
complemento para arbitrar estas transiciones.
\end{warnbox}

\subsection{Compartimentos, reticulado y dominancia}
Solo cuatro niveles son demasiado simples para una organización real. El modelo agrega una
\textbf{lista de etiquetas} (categorías) que se asignan a sujetos y objetos; un
\textbf{nivel + lista de etiquetas} forma un \textbf{compartimento}. Esto implementa el
concepto de \textbf{\emph{need-to-know}}: \emph{un general a cargo de criptografía no tiene
por qué conocer los secretos de los misiles, sin importar su nivel de acceso}.

\begin{defbox}{Relación de dominancia}
Con compartimentos, el ``mayor/menor'' se generaliza a la relación de \textbf{dominancia}.
Para niveles $(L,C)$ y $(L',C')$ —donde $C$ es el conjunto de etiquetas (categorías)—:
\[
  (L,C) \dom (L',C') \iff L' \le L \;\;\wedge\;\; C' \subseteq C.
\]
Las reglas quedan: $s$ puede \emph{leer} $o$ si $L(s)\dom L(o)$, y puede \emph{escribir} $o$
si $L(o)\dom L(s)$. Ejemplos:
\[
  (\text{TS},\{\text{nuclear}\}) \dom (\text{S},\{\text{nuclear}\}), \qquad
  (\text{S},\{\text{nuclear}\}) \dom (\text{S},\varnothing),
\]
\[
  (\text{TS},\{\text{nuclear},\text{crypto}\}) \dom (\text{TS},\{\text{nuclear}\}).
\]
\end{defbox}

Esto forma un \textbf{reticulado} (\emph{lattice}), un \textbf{conjunto de orden parcial}:
no todos los pares se pueden ordenar. Por ejemplo, $(\text{S},\{\text{nuclear}\})$ y
$(\text{S},\{\text{crypto}\})$ \emph{no} se dominan en ningún sentido (mismo nivel, pero
ninguna lista de etiquetas es subconjunto de la otra) $\Rightarrow$ están
\textbf{efectivamente separados}: ni lectura ni escritura entre ellos.

% ---- Diagrama: reticulado (recreación slide 13) ----
\begin{figure}[h]
\centering
\begin{tikzpicture}[font=\scriptsize,>=Stealth,every node/.style={align=center}]
  \node[draw,rounded corners,fill=boxBlue!40] (top) at (0,3)
       {(Top Secret,\\ \{nuclear, crypto\})};
  \node[draw,rounded corners,fill=boxBlue!25] (tn) at (-4,1.5) {(Top Secret,\\ \{nuclear\})};
  \node[draw,rounded corners,fill=boxBlue!25] (snc) at (0,1.5) {(Secret,\\ \{nuclear, crypto\})};
  \node[draw,rounded corners,fill=boxBlue!25] (tc) at (4,1.5) {(Top Secret,\\ \{crypto\})};
  \node[draw,rounded corners,fill=boxBlue!15] (sn) at (-4,-0.2) {(Secret,\\ \{nuclear\})};
  \node[draw,rounded corners,fill=boxBlue!15] (t0) at (0,-0.2) {(Top Secret,\\ \{\})};
  \node[draw,rounded corners,fill=boxBlue!15] (sc) at (4,-0.2) {(Secret,\\ \{crypto\})};
  \node[draw,rounded corners,fill=slideGray] (bot) at (0,-1.7) {(Secret, \{\})};
  \draw (top)--(tn) (top)--(snc) (top)--(tc);
  \draw (tn)--(sn) (tc)--(sc);
  \draw[dashed] (tn)--(t0) (snc)--(t0) (tc)--(t0);
  \draw (sn)--(bot) (t0)--(bot) (sc)--(bot);
\end{tikzpicture}
\caption{Reticulado (orden parcial) de compartimentos. Arriba domina a abajo; ramas sin
camino no se dominan. \textit{(Recreación de la diapositiva 13.)}}
\end{figure}

\begin{practbox}{Implementación y uso real}
La implementación es \textbf{simple}: por cada objeto se guarda como metadata un \emph{nivel}
y una \emph{lista de etiquetas}; por cada sujeto (usuario), lo mismo. Resolver una operación
es aplicar las dos reglas. Por su rigidez, \textbf{rara vez se aplica a todo un sistema}: se
reserva para lo muy crítico. Caso típico en empresas: separar por etiquetas
\emph{finanzas / producción / marketing}, con gente que accede solo a una y gente que accede a
varias. La compartimentalización tipo Bell--LaPadula aparece, p.\ ej., en sistemas que manejan
\textbf{información médica} sujetos a regulación.
\end{practbox}

%=====================================================================
\section{Modelo Biba (integridad)}
%=====================================================================

Biba desarrolló una \textbf{familia de unos 7 modelos} (en conjunto, ``modelo Biba''). Es
famoso por ser el primer modelo público que \emph{no} buscaba confidencialidad sino proteger
la \textbf{integridad} de la información. Tiene importancia en sistemas \textbf{documentales}
o de \textbf{evidencia}, donde se necesitan garantías fuertes de \emph{procedencia} y de
\emph{no adulteración}: registros contables/financieros que hay que retener, sistemas
judiciales, etc.

\begin{defbox}{Objetivo y elementos}
\textbf{Objetivo:} preservar los datos y su \textbf{integridad}. Identifica las maneras
\emph{autorizadas} en que la información puede ser alterada y \emph{quiénes} están autorizados
a hacerlo. Elementos: Sujetos, Objetos, Modos de acceso $\{r,w,\dots\}$ y
\textbf{niveles de integridad} $\il(\cdot)$, donde \emph{a mayor nivel, mayor confianza}.
\end{defbox}

Los niveles suelen interpretarse mejor por \emph{calidad/confianza} que por
\emph{secreto}: \emph{no confiable, poco confiable, confiable, muy confiable, fidedigna}.
Es estructuralmente \textbf{el inverso} de Bell--LaPadula en cuanto a reglas.

\begin{anecbox}{Una persona es tan confiable como las fuentes que usa}
La intuición del modelo: para considerar confiable a un sujeto, necesito que la información
en la que basa lo que hace o dice (lo que \emph{lee}) sea de \emph{ese nivel o más alto}. Y la
información que un sujeto \emph{genera} (escribe) será tan buena como el valor que le atribuyo
a ese sujeto: nunca mayor. De ahí salen las dos reglas.
\end{anecbox}

\subsection{Variante 1: Biba estricto}
\begin{defbox}{Reglas de Biba estricto}
\begin{enumerate}[leftmargin=1.7em,itemsep=3pt]
  \item \textbf{Escritura — \emph{Write Down}.} $s$ puede \emph{modificar} $o$ si y solo si
        \[ \il(s) \ge \il(o). \]
  \item \textbf{Lectura — \emph{Read Up}.} $s$ puede \emph{observar} $o$ si y solo si
        \[ \il(o) \ge \il(s). \]
\end{enumerate}
\textbf{Regla mnemotécnica:} \emph{Read up, Write down} (la información ``baja'' de nivel).
\end{defbox}

% ---- Diagrama Biba estricto ----
\begin{figure}[h]
\centering
\begin{tikzpicture}[font=\small,>=Stealth]
  \draw[->,line width=6pt,black!15] (7.6,2.2) -- node[right=3pt,black!55,font=\scriptsize]{Flujo Info} (7.6,-2.2);
  % Write down
  \node[circle,draw,fill=itbaBlue,minimum size=0.7cm] (sw) at (0,1.4) {};
  \node[draw,fill=black!10,minimum size=0.55cm] (ow1) at (3,2.0) {};
  \node[draw,fill=black!10,minimum size=0.55cm] (ow2) at (3,0.6) {};
  \draw[->,line width=1pt] (sw) -- node[above,font=\scriptsize]{w} (ow1);
  \draw[->,line width=1pt] (sw) -- node[below,font=\scriptsize]{w} (ow2);
  \node[font=\scriptsize\bfseries,itbaCyanDark] at (1.5,2.7) {Write Down: $\il(s)\ge\il(o)$};
  % Read up
  \node[circle,draw,fill=itbaBlue,minimum size=0.7cm] (sr) at (0,-1.6) {};
  \node[draw,fill=black!10,minimum size=0.55cm] (or1) at (3,-0.8) {};
  \node[draw,fill=black!10,minimum size=0.55cm] (or2) at (3,-2.4) {};
  \draw[<-,line width=1pt] (sr) -- node[above,font=\scriptsize]{r} (or1);
  \draw[<-,line width=1pt] (sr) -- node[below,font=\scriptsize]{r} (or2);
  \node[font=\scriptsize\bfseries,itbaCyanDark] at (1.5,-3.0) {Read Up: $\il(o)\ge\il(s)$};
\end{tikzpicture}
\caption{Biba estricto: se lee hacia arriba y se escribe hacia abajo. \textit{(Recreación de
las diapositivas de práctica.)}}
\end{figure}

\begin{practbox}{Usos del modelo estricto}
Se usa en el \textbf{curado de información}: muchos diarios lo aplican ``sin saberlo'' al
sanitizar y verificar información para generar artículos confiables. También aplica si el
\textbf{sujeto es un programa} y los \textbf{objetos son datos}: un sistema cerrado bien
clasificado evita que un programa poco robusto procese datos de mala calidad (p.\ ej.\ sin
protección contra inyección SQL). \emph{Filosofía detrás del aviso del SO:} cuando bajás un
ejecutable de Internet, el sistema te advierte ``viene de una fuente dudosa, ¿seguro querés
correrlo?''; si lo bajás del App Store, no pregunta nada. Está clasificando la confianza del
archivo por su origen: es una aplicación de la idea de Biba.
\end{practbox}

\subsection{Variante 2: Biba Ring-Policy}
\begin{defbox}{Reglas de Biba Ring-Policy}
\begin{enumerate}[leftmargin=1.7em,itemsep=3pt]
  \item \textbf{Escritura — \emph{Write Down}:} igual que el estricto, $\il(s)\ge\il(o)$.
  \item \textbf{Lectura:} \emph{cualquier} sujeto $s$ puede leer \emph{cualquier} objeto $o$
        (no se restringe la lectura).
\end{enumerate}
\end{defbox}

\subsection{Variante 3: Biba Low-Watermark}
Refleja cierto \textbf{dinamismo} (a diferencia de los estáticos). Filosofía: \emph{uno es
tan confiable como la información que usa}.

\begin{defbox}{Regla de Biba Low-Watermark}
\begin{enumerate}[leftmargin=1.7em,itemsep=3pt]
  \item \textbf{Escritura — \emph{Write Down}:} $\il(s)\ge\il(o)$.
  \item \textbf{Lectura:} cualquier sujeto puede leer cualquier objeto, \textbf{pero} al
        hacerlo su nivel de integridad se \emph{reclasifica al mínimo}:
        \[ \text{si } s \text{ lee } o:\qquad \il(s) = \min\bigl(\il(s),\, \il(o)\bigr). \]
\end{enumerate}
\end{defbox}

\begin{anecbox}{Sistemas ``adaptativos''}
Si el sistema me considera muy íntegro y en algún momento interactúo con algo de baja
integridad, mi nivel \textbf{baja automáticamente}: tendré que volver a pedir permiso o
revalidarme para recuperar el nivel anterior. Esta es la base de muchos sistemas que hoy se
venden como \emph{adaptativos / personalizados}: según las interacciones, van cambiando el
nivel de confianza del sujeto.
\end{anecbox}

%=====================================================================
\section{Modelo Clark--Wilson (integridad comercial)}
%=====================================================================

Creado por David D.~Clark y David R.~Wilson, es un modelo de \textbf{integridad} más acotado
y simple, pero introduce un concepto clave para los \textbf{sistemas operativos}. Está pensado
para escenarios donde conviven, al mismo tiempo, \textbf{información y usuarios confiables y no
confiables}.

\begin{defbox}{Objetivo y elementos}
\textbf{Objetivo:} integridad, especialmente en \textbf{entornos comerciales}. Elementos:
\begin{itemize}[leftmargin=1.4em,itemsep=2pt]
  \item \textbf{Sujetos} y \textbf{Objetos}, y estos últimos se dividen en:
    \begin{itemize}[leftmargin=1.2em,itemsep=1pt]
      \item \textbf{CDI} (\emph{Constrained Data Items}): datos restringidos / protegidos.
      \item \textbf{UDI} (\emph{Unconstrained Data Items}): datos no restringidos.
    \end{itemize}
  \item \textbf{Procedimientos} (administran accesos y permisos):
    \begin{itemize}[leftmargin=1.2em,itemsep=1pt]
      \item \textbf{TP} (\emph{Transformation Procedure}): cambia los ítems de un estado
            válido (consistente) a otro estado válido.
      \item \textbf{IVP} (\emph{Integrity Verification Procedure}): verifica que los ítems
            estén en un estado válido.
    \end{itemize}
\end{itemize}
\end{defbox}

\begin{slidebox}
\begin{itemize}[leftmargin=1.4em,itemsep=2pt]
  \item Un usuario \textbf{no puede modificar un CDI en forma directa}: debe hacerlo a través
        de un \textbf{TP}.
  \item Un usuario \textbf{sí} puede acceder a datos \textbf{UDI} sin un TP.
  \item Es responsabilidad de los TP \textbf{validar la integridad} de la información.
  \item Incorpora \textbf{\emph{separation of duty}} (separación de tareas): \emph{quien
        modifica los datos no puede ser el mismo que verifica la integridad}.
\end{itemize}
\end{slidebox}

% ---- Diagrama Clark-Wilson (recreación slide práctica usuario/CDI/UDI) ----
\begin{figure}[h]
\centering
\begin{tikzpicture}[font=\small,>=Stealth,node distance=10mm]
  \node[circle,draw,fill=itbaBlue,minimum size=0.8cm] (u) at (0,0) {};
  \node[font=\scriptsize] at (0,-0.8) {Sujeto};
  \node[draw,fill=boxBlue!50,minimum width=1.6cm,minimum height=1cm] (prog) at (4,0.4) {IVP/TP};
  \node[font=\scriptsize] at (4,-0.45) {Programa};
  \node[draw,cylinder,shape border rotate=90,fill=lvlS!40,aspect=0.35,
        minimum width=1.6cm,minimum height=1.1cm] (cdi) at (8,0.4) {CDI};
  \node[draw,cylinder,shape border rotate=90,fill=lvlU!50,aspect=0.35,
        minimum width=1.6cm,minimum height=1.1cm] (udi) at (4,-2.4) {UDI};
  \draw[<->,line width=1pt] (u) -- (prog);
  \draw[<->,line width=1pt] (prog) -- node[above,font=\scriptsize]{TP} (cdi);
  \draw[->,line width=1pt,practGreen] (u) -- (udi);
  \draw[dashed,attackRed,line width=1pt] (u) to[bend right=35] (cdi);
  \node[attackRed,font=\scriptsize\bfseries] at (4,2.0) {$\varnothing$ (prohibido el acceso directo a CDI)};
  \draw[attackRed,line width=1pt] (3.4,1.7) -- (4.6,1.7);
\end{tikzpicture}
\caption{El sujeto modifica un CDI \emph{solo} a través de un TP (con verificación IVP); a los
UDI accede directamente. \textit{(Recreación de la diapositiva de práctica.)}}
\end{figure}

\subsection{Reglas formales (clase práctica)}
Las reglas se dividen en dos grupos:

\begin{practbox}{1. Reglas de Certificación (CR) — discrecionales, las aplica el Administrador}
\begin{description}[leftmargin=2.6em,style=nextline,itemsep=2pt,font=\bfseries]
  \item[C1 (Certificación)] Los IVP aseguran que los CDI están en un estado válido.
  \item[C2 (Validez)] Un TP toma un estado válido y lo lleva a otro estado válido,
        sobre un conjunto certificado $\{TP_i, (CDI_a, CDI_b, \dots)\}$.
  \item[C3 (Separación de tareas)] Las \emph{relaciones certificadas} que asocian un conjunto
        de CDI con un TP cumplen la separación de tareas.
  \item[C4 (Journal certification)] En un CDI tipo bitácora, los TP guardan información para
        reconstruir la operación.
  \item[C5 (Validación de UDI)] Los TP validan los UDI: o convierten el UDI en un CDI válido,
        o el UDI es rechazado.
\end{description}
\end{practbox}

\begin{practbox}{2. Reglas de Aplicación (ER) — mandatorias, las aplica el Sistema}
\begin{description}[leftmargin=2.6em,style=nextline,itemsep=2pt,font=\bfseries]
  \item[E1 (Aplicación de validez)] El sistema mantiene la lista de relaciones certificadas:
        $TP_f$ opera sobre $CDI_o \iff (f,o)\in C$.
  \item[E2 (Aplicación de separación)] El sistema mantiene la lista de relaciones permitidas
        $(UserID, TP_i, \{CDI_a,\dots\})$.
  \item[E3 (Identificación de usuarios)] El sistema debe \emph{autenticar} a quien intenta
        ejecutar el TP.
  \item[E4 (Cambios)] Las autorizaciones solo las puede cambiar el administrador
        (\emph{security officer}).
\end{description}
\end{practbox}

% ---- Diagrama Clark-Wilson formal (recreación slide formal) ----
\begin{figure}[h]
\centering
\begin{tikzpicture}[font=\scriptsize,>=Stealth,node distance=8mm]
  \node[font=\bfseries] (users) at (0,2.4) {USERS};
  \node[draw,fill=boxBlue!50,minimum width=1.3cm,minimum height=0.9cm] (tp) at (0,0.8) {TP};
  \node[draw,signal,signal to=east,fill=lvlS!40] (udi) at (-3.2,0.8) {UDI};
  \node[draw,starburst,fill=lvlS!30,inner sep=1pt,minimum size=0.9cm] (ivp) at (0,-1.0) {IVP};
  \node[draw,circle,fill=lvlU!40,minimum size=0.8cm] (cdia) at (3.4,2.0) {$CDI_a$};
  \node[draw,circle,fill=lvlU!40,minimum size=0.8cm] (cdib) at (3.4,0.8) {$CDI_b$};
  \node[draw,circle,fill=lvlU!40,minimum size=0.8cm] (log) at (3.4,-0.5) {Log};
  \node[draw,circle,fill=lvlU!40,minimum size=0.8cm] (cdin) at (3.4,-1.7) {$CDI_n$};
  \draw[->] (users) -- node[left,font=\tiny]{C3} (tp);
  \draw[->] (udi) -- node[above,font=\tiny]{C5} (tp);
  \draw[->] (tp) -- node[above,font=\tiny]{C2} (cdia);
  \draw[->] (tp) -- node[above,font=\tiny]{C2} (cdib);
  \draw[->] (tp) -- node[above,font=\tiny]{C4} (log);
  \draw[<-] (ivp) -- node[below,font=\tiny]{C1} (cdin);
  \node[font=\scriptsize\itshape,black!60] at (3.4,-2.6) {System in valid state};
\end{tikzpicture}
\caption{Reglas de certificación: el TP lleva los CDI de un estado válido a otro, el IVP los
verifica, los UDI se validan y todo queda registrado. \textit{(Recreación de la diapositiva
``Clark--Wilson (Formal)''.)}}
\end{figure}

\begin{anecbox}{De los bancos al kernel: separation of duty en acción}
El modelo se aplica también a sistemas \emph{no} informáticos. Es la base del
\textbf{principio de separación de privilegios}. En un banco, al abrir una cuenta de empresa
se arma un \emph{registro de firmas}: por debajo de cierto monto bastan ciertas firmas; por
encima se exige una firma más. Para procesar una transacción siempre hay \emph{dos partes}: una
\textbf{verifica} las firmas y otra, \textbf{distinta}, \textbf{ejecuta} la transacción
—aunque todo sea automático—.\medskip

\noindent En sistemas operativos, esta idea derivó en la separación
\textbf{\emph{user space} / \emph{kernel space}}: el \emph{kernel space} es un ambiente de alta
confianza (puede saltarse todos los controles: acceso directo al disco, etc.); el resto corre
como \emph{no confiable} en \emph{user space}. ¿Cómo se cruza de uno a otro? Con las
\textbf{\emph{syscalls}} (llamadas al sistema): el usuario arma un \emph{dataset} describiendo
qué quiere (``abrir tal archivo para lectura, con tal \emph{file descriptor}\dots'') y dispara
la syscall. Un proceso \textbf{transforma} (TP) esa información llevándola a la zona del
kernel, otro \textbf{verifica} su integridad (IVP) (que el \emph{file descriptor} sea un
número positivo, etc.) y, recién entonces, un programa privilegiado ejecuta. Es una aplicación
directa de Clark--Wilson para proteger la integridad del sistema operativo.
\end{anecbox}

%=====================================================================
\section{Más allá de los modelos teóricos}
%=====================================================================

Los tres modelos vistos son \textbf{teóricos}: en la práctica \emph{no se aplican
directamente} a un sistema completo, pero son la base de muchas implementaciones. Es raro
encontrar el nombre explícito de Bell--LaPadula / Biba / Clark--Wilson, porque mucha gente
aplica la idea sin saber el nombre. Aun así, conviene saber dónde aparecen:

\begin{center}
\renewcommand{\arraystretch}{1.35}
\begin{tabular}{>{\bfseries}p{3.4cm} p{9.6cm}}
\toprule
Modelo & ¿Qué sistemas se basan en él? \\
\midrule
Bell--LaPadula & La organización por \emph{compartimentos} aparece en sistemas que manejan
\textbf{información médica} (sujetos a regulación).\\
Biba & Sistemas de \emph{clasificación y archivo} de información para \textbf{medios de prensa,
legajos judiciales o declaraciones fiscales}.\\
Clark--Wilson & El concepto de \textbf{\emph{kernel space} / \emph{user space}} de muchos
\textbf{sistemas operativos} (p.\ ej.\ Unix).\\
\bottomrule
\end{tabular}
\end{center}

\begin{practbox}{Lectura recomendada sobre los modelos (papers originales)}
\begin{itemize}[leftmargin=1.4em,itemsep=1pt]
  \item \texttt{Bell76.pdf} — documento original del modelo Bell--LaPadula.
  \item \texttt{Biba75.pdf} — documento original del modelo Biba.
  \item \texttt{A\_Comparison\_of\_commercial\_and\_military.pdf} — documento original del
        modelo Clark--Wilson.
\end{itemize}
\end{practbox}

%=====================================================================
\section{Mecanismos de control de acceso: MAC vs.\ DAC}
%=====================================================================

\begin{slidebox}
\centering El \textbf{control de acceso} es un mecanismo de seguridad que \textbf{regula
quién (sujeto) puede acceder a qué (objeto) y bajo qué condiciones}.
\end{slidebox}

\medskip
Los modelos suelen ser muy restrictivos, así que rara vez cubren todo un sistema; en el resto
necesitamos manejar el control de acceso con dos grandes familias:

\begin{center}
\renewcommand{\arraystretch}{1.3}
\begin{tabular}{>{\raggedright\arraybackslash}p{6.4cm} >{\raggedright\arraybackslash}p{6.4cm}}
\toprule
\highlight{macGreen}{\bfseries MAC — Mandatorio} & \highlight{dacYellow}{\bfseries DAC — Discrecional}\\
\emph{Basado en reglas} & \emph{Basado en identidad}\\
\midrule
Control que se aplica en forma \textbf{obligatoria a todos los accesos} de los sujetos a los
objetos. Reglas universales (p.\ ej.\ Bell--LaPadula Lattice). Las políticas se definen
\textbf{para el sistema}. &
La asignación y modificación de permisos queda \textbf{a discreción} del sujeto responsable
del objeto (o de un administrador). Las políticas se definen \textbf{para el usuario}.\\[3pt]
Aplica sobre: todos los accesos, los cambios en las reglas/atributos, y el intercambio de
información entre sujetos. &
Típico de \textbf{repositorios donde los sujetos crean los objetos} (file systems, object
storage). Permite distribuir la responsabilidad.\\
\bottomrule
\end{tabular}
\end{center}

\begin{anecbox}{¿Por qué existe el DAC?}
Los sistemas mandatorios son demasiado rígidos para sistemas de \emph{propósito general} (un
sistema operativo, una base de datos, Google Workspace o Microsoft~365). En el DAC las
propiedades de seguridad \emph{no} las da el sistema, sino el administrador según cómo lo
configure: es una forma elegante de cubrir múltiples escenarios y de no tener que dar garantías
fuertes. La contracara: las garantías dependen de \emph{cómo} se configura.
\end{anecbox}

%=====================================================================
\section{Control discrecional por extensión}
%=====================================================================

Dentro del DAC, los accesos se pueden definir \textbf{por extensión} (enumerando) o
\textbf{por comprensión} (mediante reglas/propiedades). Empecemos por extensión.

\subsection{Matriz de control de accesos}
Representación matricial: en un eje los \textbf{sujetos} (filas), en el otro los
\textbf{objetos} (columnas); en cada intersección, las \textbf{operaciones permitidas} para
esa combinación.

% ---- Diagrama matriz / ACL / capability (recreación slide práctica) ----
\begin{figure}[h]
\centering
\begin{tikzpicture}[font=\small]
  \matrix (m) [matrix of nodes, nodes={draw,minimum width=1.7cm,minimum height=0.7cm,anchor=center},
               row 1/.style={nodes={fill=lvlS!50,font=\bfseries}},
               column 1/.style={nodes={fill=attackRed!15,font=\bfseries}},
               column sep=-\pgflinewidth, row sep=-\pgflinewidth]
  {
     & Obj1 & Obj2 & $\cdots$ & ObjN \\
    Usuario1 & r & rw & & rwx \\
    Usuario2 & rx & r & & rx \\
    $\cdots$ & & & & \\
  };
  % ACL = columna
  \draw[decorate,decoration={brace,amplitude=5pt},attackRed,line width=1pt]
        (m-1-2.north west) -- (m-1-5.north east)
        node[midway,above=5pt,font=\scriptsize\bfseries]{ACL (columnas)};
  % capability list = fila
  \draw[decorate,decoration={brace,amplitude=5pt,mirror},practGreen,line width=1pt]
        (m-2-1.north west) -- (m-4-1.south west)
        node[midway,left=5pt,font=\scriptsize\bfseries,align=center,rotate=90,anchor=south]{Capability lists (filas)};
\end{tikzpicture}
\caption{Matriz de accesos. Sus \emph{columnas} son las \textbf{ACL} (por objeto); sus
\emph{filas} son las \textbf{capability lists} (por sujeto). \textit{(Recreación de las
diapositivas de práctica.)}}
\end{figure}

\begin{warnbox}{Inconvenientes de la matriz}
Es simple de entender pero, como implementación directa, tiene problemas serios. Sirve más
como \emph{concepto ilustrativo}:
\begin{itemize}[leftmargin=1.4em,itemsep=1pt]
  \item \textbf{Escalabilidad}: con $100\,000$ usuarios y $1\,000\,000$ de objetos crece de
        forma violenta.
  \item \textbf{Eficiencia}: la búsqueda puede consumir mucho tiempo y recursos.
  \item \textbf{Mantenimiento}: agregar un sujeto obliga a definir su relación con \emph{todos}
        los objetos (y viceversa). Propenso a errores.
  \item \textbf{Almacenamiento} y \textbf{limitaciones dinámicas}: no se adapta bien a entornos
        donde cambian permisos, sujetos u objetos.
\end{itemize}
\textbf{Clave (consulta de Joaquín):} para la \emph{computadora} tener una matriz gigante no
es problema; el problema es el \emph{pobre humano} que con su discreción tiene que poner los
valores. \emph{Eso} es lo que no escala.
\end{warnbox}

\subsection{Listas de control de acceso (ACL)}
En un sistema grande los permisos aparecen como ``islas'': la matriz queda \textbf{dispersa}
(casi todo vacío). Una forma eficiente de representarla es, por cada objeto, enumerar solo los
pares \emph{(sujeto, permisos)} que tienen sentido. Eso son las \textbf{ACL}: son las
\emph{columnas} de la matriz.

\begin{defbox}{Definición formal de ACL}
\[
  \mathrm{ACL}(o) = \{\, (s_i, p_i) \mid s_i \in \Sset \ \wedge\ p_i \subseteq \Rset \,\}
\]
con $\Sset=$ sujetos, $\Oset=$ objetos, $\Rset=$ permisos. Extensiones habituales:
\begin{itemize}[leftmargin=1.4em,itemsep=1pt]
  \item \textbf{Denegar por defecto}: si $s\notin\mathrm{ACL}(o)\Rightarrow s$ no tiene
        permiso sobre $o$.
  \item \textbf{Grupos}: $G=\{s_1,s_2,\dots,s_i \mid s_i\in\Sset\}$, p.\ ej.\
        $\mathrm{ACL}(o)=\{(s_1,r),(G,rw),\dots\}$. Permite tratar a muchos sujetos como una
        sola entrada y aplicar reglas a nivel grupo.
  \item \textbf{Wildcards}: $\mathrm{ACL}(o)=\{(s_1,r),(s_2,\ast),\dots\}$.
  \item \textbf{Permisos por defecto} (\emph{defaults}) y \textbf{patrones} para identificar
        objetos: \texttt{/users/*}, \texttt{/home/**/*.tmp}.
\end{itemize}
\end{defbox}

\subsubsection{Conflictos en ACL}
Si \textbf{más de una entrada} de la ACL otorga permisos a un sujeto $s$, hay que resolver:
\begin{enumerate}[leftmargin=1.7em,itemsep=2pt]
  \item \textbf{Permitir} el acceso si \emph{alguna} entrada lo da.
  \item \textbf{Grant All}: denegar si \emph{alguna} entrada lo deniega (exige que
        \emph{todas} den acceso).
  \item \textbf{First Rule / First Match}: aplicar la \emph{primera} regla encontrada.
\end{enumerate}

\begin{practbox}{Ejercicio resuelto — conflictos de ACL}
Sea $s_1 \in G$ y
\[
  \mathrm{ACL}(o) = \{\, (G,\,rw),\ (s_1,\,w),\ (\{s_1,s_2\},\,wx),\ \dots \,\}.
\]
¿Qué puede hacer $s_1$ según cada política?
\begin{center}
\renewcommand{\arraystretch}{1.25}
\begin{tabular}{>{\bfseries}p{3.6cm} p{4.2cm} p{4.6cm}}
\toprule
Política & Entradas que aplican & Resultado para $s_1$\\
\midrule
Permitir (unión) & $rw,\ w,\ wx$ & \highlight{slideGray}{$r,w,x$} (la unión)\\
Grant All (intersección) & $rw \cap w \cap wx$ & \highlight{slideGray}{$w$}\\
First Rule & primera: $(G,rw)$ & \highlight{slideGray}{$r$ y $w$}\\
\bottomrule
\end{tabular}
\end{center}
\end{practbox}

\subsubsection{Derechos por defecto (clase práctica)}
Otras reglas aplicables cuando hay una entrada específica y un \emph{default}:
\begin{enumerate}[leftmargin=1.7em,itemsep=2pt]
  \item \textbf{Override / Best match}: si el sujeto tiene una entrada específica, usarla;
        si no, usar el \emph{default}.
  \item \textbf{Augment}: partir del \emph{default} y \textbf{agregarle} las entradas
        explícitas.
\end{enumerate}

\begin{practbox}{Ejercicio resuelto — override vs.\ augment}
Sea $\mathrm{ACL}(o) = \{\, (s_1,\,x),\ (\ast,\,r),\ \dots \,\}$ (el \emph{default} ``$\ast$''
da $r$ a todos). ¿Qué puede hacer $s_1$?
\begin{itemize}[leftmargin=1.4em,itemsep=1pt]
  \item \textbf{Override}: usa la entrada específica $\Rightarrow$ puede \highlight{slideGray}{$x$}.
  \item \textbf{Augment}: default $+$ específica $\Rightarrow$ puede
        \highlight{slideGray}{$r$} (porque todos pueden) \highlight{slideGray}{$x$}.
\end{itemize}
\end{practbox}

\subsection{Usuarios privilegiados, transferencias y revocaciones}
\begin{practbox}{Operaciones sobre las ACL (clase práctica)}
\begin{description}[leftmargin=2.2em,style=nextline,itemsep=3pt,font=\bfseries]
  \item[Usuarios privilegiados]
    \begin{enumerate}[leftmargin=1.2em,itemsep=0pt]
      \item \emph{Root} o administrador.
      \item \emph{Owner}: puede modificar la ACL del objeto sobre el que tiene el derecho
            ``\texttt{own}''.
    \end{enumerate}
  \item[Transferencia] Cambia el sujeto de una entrada: el sujeto\,1 \emph{deja} de tener los
        derechos, que pasan al sujeto\,2.
  \item[Delegación] El sujeto\,1 le \emph{delega} derechos al sujeto\,2 \textbf{sin perderlos}.
  \item[Revocación] El \emph{owner} quita la entrada o quita derechos. Si hubo delegación,
        se produce \textbf{revocación en cascada}.
\end{description}
\end{practbox}

%=====================================================================
\section{Control discrecional por comprensión}
%=====================================================================

Para sistemas grandes no alcanza con la representación dispersa: en cada ``isla'' hay grupos
de usuarios y objetos con casi la misma distribución de accesos. La solución es definir el
control \textbf{por comprensión} (mediante \emph{reglas}/propiedades) en vez de por extensión,
guardando reglas compactas que podrían reconstruir la matriz \emph{bajo demanda}. El criterio
de agrupación define el mecanismo:

\begin{center}
\renewcommand{\arraystretch}{1.3}
\begin{tabular}{>{\bfseries}p{2.6cm} p{2.8cm} p{7.4cm}}
\toprule
Sigla & Por\dots & Idea \\
\midrule
RBAC & roles & \emph{Role-Based Access Control}.\\
ABAC & atributos & \emph{Attribute-Based Access Control}.\\
RuBAC & reglas & \emph{Rule-Based Access Control}.\\
\bottomrule
\end{tabular}
\end{center}

\noindent No hay uno mejor: según la situación conviene uno u otro.

\subsection{RBAC — control basado en roles}
Estereotipa \textbf{roles} (más abarcativo que un grupo): genera ``plantillas'' que agrupan
permisos. Ejemplos de roles en una empresa: \emph{account manager}, persona de ventas,
administrador de IT, administrador de operaciones, usuario común; o
\emph{gestor de pacientes}, \emph{analista de compras}, \emph{DBA}.
\begin{itemize}[leftmargin=1.4em,itemsep=1pt]
  \item Los \textbf{permisos se asignan a los roles}; los \textbf{roles se asignan a los
        sujetos} (o a grupos).
  \item Considera \textbf{jerarquías de roles}. En sistemas tradicionales un usuario tiene un
        solo rol, pero las aplicaciones complejas permiten \emph{varios}.
  \item Una vez definidos los roles, la \textbf{gestión diaria es simple}: a un usuario nuevo
        se le dan los roles correctos.
\end{itemize}

\begin{anecbox}{¿RBAC es MAC o DAC?}
Algunos libros lo describen como \textbf{mandatorio} (porque define ``reglas de juego
universales''), pero es una visión sesgada. Hay una \textbf{discrecionalidad fuerte}: alguien
tiene que \emph{asignar} usuarios a roles, y las garantías dependen de cómo se haga esa
asignación —característica que define a lo discrecional—. Además, cuando un usuario puede tener
\textbf{varios roles}, aparece un problema interesante: \emph{múltiples sistemas de control de
acceso interactuando sobre la misma operación} (p.\ ej.\ un esquema de roles más una zona
protegida por Bell--LaPadula), lo que obliga a resolver conflictos.
\end{anecbox}

\subsection{ABAC — control basado en atributos}
Las reglas se basan en \textbf{atributos}: del sujeto (rol, edad\dots), del objeto (tipo,
clasificación de seguridad\dots), de la acción y del contexto (tiempo, ubicación\dots).

\subsection{RuBAC — control basado en reglas}
Muy parecido a roles, pero \textbf{sin asignación de sujetos}. Una regla en su forma más
básica contiene: \textbf{Principal}, \textbf{Objeto}, \textbf{Tipo de acceso} y si está
\textbf{permitido o bloqueado}. Típicamente hay una \textbf{lista ordenada de reglas} y
``aplica la primera que ocurra'' (u otra política de resolución).

%=====================================================================
\section{Resolución de múltiples reglas}
%=====================================================================

Cuando varias reglas aplican de forma \textbf{contradictoria}, se usa alguna política de
resolución. \textbf{El mismo juego de reglas resuelve muy distinto según cómo se manejen los
conflictos} (clave si alguna vez configuran firewalls):

\begin{defbox}{Políticas de resolución}
\begin{itemize}[leftmargin=1.4em,itemsep=2pt]
  \item \textbf{First match}: las reglas tienen un orden; aplica la \emph{primera} que hace
        \emph{match}.
  \item \textbf{Best match}: aplica la regla \emph{más específica} (las reglas más específicas
        priman sobre las generales).
  \item \textbf{Deny over Allow}: cualquier regla que \emph{deniega} explícitamente tiene
        precedencia sobre las que permiten (deben permitir \emph{todas} las que apliquen).
\end{itemize}
\end{defbox}

\begin{practbox}{Ejemplo: reglas de Firewall}
Un \emph{firewall} arbitra el tráfico entre dos redes. En su forma más simple es un sistema
de reglas ordenadas. Considérese:
\begin{center}
\ttfamily\footnotesize
\renewcommand{\arraystretch}{1.1}
\begin{tabular}{r l l l l l}
\toprule
\# & Proto & Source IP & Dest IP & Port & Action\\
\midrule
1 & TCP & 10.1.1.1 & 20.1.1.1 & 80 & Accept\\
2 & TCP & 10.1.1.2 & 20.1.1.1 & 80 & Deny\\
3 & TCP & 10.1.1.0/24 & 20.1.1.1 & 80 & Deny\\
4 & TCP & 10.1.1.3 & 20.1.1.1 & 80 & Accept\\
5 & TCP & 10.2.2.0/24 & 20.2.2.5 & 80 & Deny\\
6 & TCP & 10.2.2.5 & 20.2.2.0/24 & 80 & Deny\\
7 & TCP & 10.3.3.0/24 & 20.3.3.9 & 80 & Accept\\
8 & TCP & 10.3.3.9 & 20.3.3.0/24 & 80 & Deny\\
9 & IP  & 0.0.0.0/0 & 0.0.0.0/0 & 0--65535 & Deny\\
\bottomrule
\end{tabular}
\end{center}
\medskip
\textbf{Caso A — paquete \texttt{10.1.1.1 $\Rightarrow$ 20.1.1.1:80}} (hacen \emph{match} las
reglas 1 y 3):
\begin{itemize}[leftmargin=1.4em,itemsep=1pt]
  \item \emph{First match} $\Rightarrow$ regla 1 $\Rightarrow$ \textbf{Accept}.
  \item \emph{Deny over Allow} $\Rightarrow$ la 3 deniega $\Rightarrow$ \textbf{Deny}.
\end{itemize}
\textbf{Caso B — paquete \texttt{10.3.3.9 $\Rightarrow$ 20.3.3.9:80}} (hacen \emph{match} las
reglas 7 —red \texttt{/24}— y 8 —host específico—):
\begin{itemize}[leftmargin=1.4em,itemsep=1pt]
  \item \emph{Best match} $\Rightarrow$ la 8 es más específica $\Rightarrow$ \textbf{Deny}.
\end{itemize}
\textbf{Observación:} en \emph{First match}, estas reglas dicen ``no permitas comunicación
hacia la red 10.1.1.x \emph{salvo} desde el host 10.1.1.1''; para expresar lo mismo con
\emph{Best match} hay que escribirlas distinto (``permití a toda la red \emph{salvo} a tal
host''). La regla 9 final ($0.0.0.0/0$ Deny) implementa el \emph{denegar por defecto}.
\end{practbox}

%=====================================================================
\section{Reglas basadas en contexto}
%=====================================================================

Permiten construir reglas que incluyen \textbf{información de contexto}, no solo la entrada/
salida de la operación. El \textbf{contexto} lo forman todos los datos asociados al sujeto, el
objeto, el momento o el lugar del acceso —elementos que cambian \emph{independientemente} de
las reglas—:
\begin{itemize}[leftmargin=1.4em,itemsep=1pt]
  \item Momento del día; ubicación geográfica del origen de la conexión.
  \item Departamento al que pertenece la persona; si el usuario estaba de vacaciones.
  \item Nivel de confiabilidad de la máquina (p.\ ej.\ del DBA).
\end{itemize}

Permiten detectar \textbf{situaciones anómalas}. Ejemplos:
\begin{itemize}[leftmargin=1.4em,itemsep=1pt]
  \item Un usuario se loguea desde China y desde Argentina con 5 minutos de diferencia.
  \item Una usuaria que figura de licencia por maternidad se loguea en una base de datos del
        \emph{datacenter}.
  \item Un usuario inicia dos VPNs en menos de 5 minutos con dos IPs distintas.
\end{itemize}

\begin{anecbox}{El nicho: prevención de fraude}
Estas reglas son raras (muy sofisticadas) pero tienen un nicho claro: la \textbf{prevención de
fraude} en e-commerce y sistemas financieros. Distinción importante:
\begin{itemize}[leftmargin=1.4em,itemsep=1pt]
  \item Un \textbf{problema de seguridad} es un abuso que se \emph{saltea} las reglas del
        sistema.
  \item Un \textbf{fraude} \emph{sigue} las reglas del sistema, pero hace algo que el sistema
        no esperaba (p.\ ej.\ pagar con una tarjeta de crédito robada: el pago ``sale bien'',
        pero el dueño reclamará y se anulará).
\end{itemize}
El control anti-fraude estima la \emph{probabilidad} de que una operación sea fraudulenta
(en el login, el checkout o el panel de administración) y la deniega ante indicios.
\end{anecbox}

\subsection{Zero-Trust y Open Policy Agent (OPA)}
La lógica de estas reglas suele requerir \textbf{programación}, y muchos sistemas no se pueden
modificar. En el modelo \textbf{Zero-Trust}, la evaluación de contexto se define como una
\textbf{entidad externa} al \emph{Policy Engine}: un tercero valida cada operación.

\begin{defbox}{Open Policy Agent (OPA)}
\textbf{OPA} es un \emph{framework} estándar que unifica, mediante un lenguaje programable, la
implementación de \textbf{políticas de decisión}, \textbf{desacoplando} la política de control
de acceso del sistema controlado (\emph{policy decoupling}). Las políticas se escriben en un
lenguaje llamado \textbf{Rego}. El servicio que necesita una decisión \emph{delega} en OPA:
le envía una \emph{query} (cualquier valor JSON) y recibe una \emph{decision} (otro JSON); OPA
evalúa contra su \emph{Policy} (Rego) y sus \emph{Data} contextuales (JSON, que se actualizan
o se proveen \emph{on-demand}).
\end{defbox}

% ---- Diagrama OPA (recreación slide OPA) ----
\begin{figure}[h]
\centering
\begin{tikzpicture}[font=\scriptsize,>=Stealth,node distance=10mm]
  \node[draw,rounded corners,minimum width=2.6cm,minimum height=0.8cm,fill=boxBlue!40] (svc) at (0,1.6) {Service};
  \node[above=4mm of svc,font=\scriptsize] (req) {Request, Event\dots};
  \node[draw,rounded corners,minimum width=1.6cm,minimum height=0.8cm,fill=slideGray] (opa) at (0,-0.4) {OPA};
  \node[draw,fill=lvlS!50,minimum width=1.3cm,minimum height=0.8cm,align=center] (pol) at (-2.4,-2.2) {Policy\\(Rego)};
  \node[draw,cylinder,shape border rotate=90,fill=black!20,aspect=0.3,minimum width=1.2cm,align=center] (data) at (2.4,-2.2) {Data\\(JSON)};
  \draw[->] (req)--(svc);
  \draw[->] (svc.south) to[bend right=25] node[left,font=\tiny]{Query} (opa.north);
  \draw[->] (opa.north) to[bend right=25] node[right,font=\tiny]{Decision} (svc.south);
  \draw[-] (opa)--(pol); \draw[-] (opa)--(data);
\end{tikzpicture}
\caption{\emph{Policy decoupling}: el servicio delega la decisión de acceso en OPA, que evalúa
política (Rego) y datos de contexto. \textit{(Recreación de la diapositiva ``Open Policy
Agent''.)}}
\end{figure}

\begin{slidebox}
\textbf{Ejemplo de Rego} (\emph{admission control} en Kubernetes: solo imágenes del registro
\texttt{hooli.com}):
\begin{lstlisting}[style=code]
package kubernetes.admission
import rego.v1

deny contains reason if {
    some container
    input_containers[container]
    not startswith(container.image, "hooli.com/")
    reason := "container image refers to illegal registry (must be hooli.com)"
}

input_containers contains container if {
    container := input.request.object.spec.containers[_]
}
\end{lstlisting}
\end{slidebox}

%=====================================================================
\section{Ejemplos reales de control de acceso}
%=====================================================================

\subsection{Unix / Linux file system}
En el paradigma Unix, los \textbf{usuarios} son sujetos y \textbf{todo es un archivo}
(monitor, memoria, disco, mouse, teclado\dots). Las acciones se clasifican en tres:
\textbf{lectura ($r$)}, \textbf{escritura ($w$)} y \textbf{ejecución ($x$)}. Cada usuario
pertenece a uno o más \textbf{grupos}, y todo archivo tiene un \emph{usuario} y un
\emph{grupo} dueños. Los permisos se asignan en \textbf{tres secciones}: usuario, grupo y
resto (\emph{others}).

% ---- Diagrama permisos rwx (recreación slide Unix) ----
\begin{figure}[h]
\centering
\begin{tikzpicture}[font=\small,>=Stealth]
  \node[font=\Large\ttfamily\bfseries] (perm) at (0,0) {-\,rwx\,rwx\,rwx};
  \node[font=\scriptsize,align=center] (ft) at (-2.7,-1.6) {File\\ Type};
  \node[font=\scriptsize] (us) at (-1.4,-1.2) {User};
  \node[font=\scriptsize] (gr) at (-0.1,-1.2) {Group};
  \node[font=\scriptsize] (ot) at (1.3,-1.2) {Others};
  \draw[->,black!60] (ft) -- (-2.55,-0.3);
  \draw[->,itbaCyanDark] (us) -- (-1.4,-0.3);
  \draw[->,practGreen] (gr) -- (-0.1,-0.3);
  \draw[->,attackRed] (ot) -- (1.25,-0.3);
  \node[font=\scriptsize,anchor=west] at (3.0,0.5) {$r=$ readable};
  \node[font=\scriptsize,anchor=west] at (3.0,0.1) {$w=$ writeable};
  \node[font=\scriptsize,anchor=west] at (3.0,-0.3) {$x=$ executable};
  \node[font=\scriptsize,anchor=west] at (3.0,-0.7) {$-=$ denied};
\end{tikzpicture}
\caption{Permisos Unix: tipo de archivo + tres ternas \texttt{rwx} (usuario, grupo, otros).
\textit{(Recreación de la diapositiva ``Unix File System''.)}}
\end{figure}

\begin{practbox}{\texttt{umask}: herencia con máscara}
Al crear un archivo, este hereda inicialmente los atributos del directorio. El comando
\textbf{\texttt{umask}} (variable de entorno) le indica al SO que \emph{saque} permisos según
una máscara: p.\ ej.\ ``que los archivos nuevos nunca sean ejecutables'' o ``que solo los pueda
ver yo'', sin importar el directorio.
\end{practbox}

\subsection{SELinux (Security-Enhanced Linux)}
Conjunto de modificaciones al \emph{kernel} de Linux que agrega un \textbf{control de seguridad
adicional al DAC} (\emph{hardening} de servidores). Es un conjunto de reglas que siempre
responden la misma pregunta: \emph{¿puede este sujeto acceder a este objeto?}
\begin{itemize}[leftmargin=1.4em,itemsep=1pt]
  \item Cada proceso tiene un \textbf{contexto SELinux}; todos los recursos accesibles también
        tienen contexto: \texttt{/var/tmp} $\to$ \texttt{t\_var\_tmp};
        \texttt{Port 80} $\to$ \texttt{t\_port\_www};
        \texttt{/var/www/html} $\to$ \texttt{t\_web\_html}.
  \item Usa \textbf{macros} que expanden a conjuntos de reglas (limitan los \emph{syscalls}
        entre combinaciones sujeto–objeto).
  \item Potente: \textbf{dos procesos corriendo como \texttt{root} (UID 0)} pueden tener
        restringido el acceso a los recursos del otro gracias a SELinux (p.\ ej.\ Apache
        \texttt{httpd\_t} y MariaDB \texttt{mysqld\_t} aislados entre sí).
\end{itemize}

\subsection{Windows file system}
El dueño de cada archivo/directorio define los permisos (\textbf{DAC}); los sujetos pueden ser
cuentas o grupos. Tiene un nivel \textbf{mucho más granular} (11--13 operaciones por archivo) y
—a diferencia de Unix— \textbf{herencia de permisos} que \emph{persiste} tras la creación.
\begin{itemize}[leftmargin=1.4em,itemsep=1pt]
  \item Cada entrada puede \emph{Allow}, \emph{Deny} o dejar ambos en blanco (vale el heredado).
  \item \textbf{Deny tiene precedencia sobre Allow}; las \textbf{reglas locales} priman sobre
        las \textbf{heredadas}.
  \item Por su complejidad, Windows ofrece un \textbf{simulador de permisos efectivos} que
        resuelve todas las reglas y conflictos. (Históricamente, resolver permisos exigía leer
        muchos lugares del disco $\Rightarrow$ ``microcortes'' al acceder a archivos.)
\end{itemize}

\subsection{PostgreSQL: roles y policies}
\begin{itemize}[leftmargin=1.4em,itemsep=1pt]
  \item Usa un modelo \textbf{RBAC}: un rol puede asociarse a un usuario, a un grupo o a otro
        rol. El atributo \texttt{LOGIN} permite iniciar sesión. Todo objeto creado por un rol
        tiene a ese rol como \emph{owner} (otro caso de \textbf{DAC}): solo \emph{superusers} y
        el \emph{owner} acceden.
  \item Las \textbf{policies} crean políticas de acceso/modificación de \emph{registros} dentro
        de las tablas: permiten un \emph{need-to-know} a nivel de fila (\emph{Row Level
        Security}).
\end{itemize}
\begin{slidebox}
\begin{lstlisting}[style=code]
ALTER TABLE clients ENABLE ROW LEVEL SECURITY;
CREATE POLICY salesperson_clients ON clients
  USING (salesperson_id::text = current_user);
\end{lstlisting}
\end{slidebox}
\begin{anecbox}{Vendedores celosos de sus contactos}
Caso típico: un equipo de ventas donde cada vendedor (va a comisión) solo debe ver
\emph{sus} clientes. Con una \emph{policy} sobre la tabla \texttt{clients} que compara
\texttt{salesperson\_id} con el \texttt{current\_user}, una consulta sobre toda la tabla
devuelve solo las filas propias (es como agregar un filtro automático). \emph{Salvedad:} las
prácticas modernas tienden a sacar la gestión de usuarios fuera de la base de datos (la base
es buena gestionando datos, no logins), así que estos \emph{features} se ven más en el modelo
``usuarios de la app $=$ usuarios de la base''.
\end{anecbox}

\subsection{Squid Web Proxy}
Un \emph{web proxy} controla por dónde navegan los usuarios (y acelera/cachea). Squid
implementa un sistema de \textbf{reglas ordenadas} (\textbf{MAC}; ``mal llamado ACL'') para
permitir o bloquear sitios.
\begin{itemize}[leftmargin=1.4em,itemsep=1pt]
  \item Filtra por dominio, protocolo, URL, \emph{headers} del requerimiento y
        \textbf{listas negras}.
  \item En la facultad, los sitios de \emph{streaming} se bloquean con esta tecnología.
  \item \textbf{Blacklists dinámicas}: iniciativas que recolectan IPs de origen de ataques,
        metodologías y \emph{malware} usado; Squid puede consultarlas para bloquear accesos.
\end{itemize}

\subsection{AWS (IAM)}
AWS (más de 1000 servicios) combina \textbf{dos mecanismos}:
\begin{center}
\renewcommand{\arraystretch}{1.3}
\begin{tabular}{>{\raggedright\arraybackslash}p{6.2cm} >{\raggedright\arraybackslash}p{6.2cm}}
\toprule
\textbf{Identity-based policies} & \textbf{Resource-based policies}\\
\highlight{macGreen}{Asociadas al \emph{Principal} (MAC)} &
\highlight{dacYellow}{Asociadas al \emph{Recurso} (DAC)}\\
\midrule
Reglas asociadas a la cuenta/usuario; las define el \textbf{administrador}. Cubren ``lo
grueso'' de la infraestructura (subredes, firewalls). &
Quien \emph{crea} el recurso define quién accede (p.\ ej.\ una \emph{S3 bucket policy}).
Más ágil para los equipos de desarrollo.\\
\bottomrule
\end{tabular}
\end{center}

\subsection{API gateway}
Servicio que se despliega \emph{delante} de las APIs para una capa extra de seguridad.
\begin{itemize}[leftmargin=1.4em,itemsep=1pt]
  \item Controles \textbf{basados en reglas (MAC)} por URL, dominio, \emph{headers}, método.
  \item Reglas \textbf{por contexto}: máxima cantidad de requerimientos por segundo, ancho de
        banda; si se superan los umbrales, bloquea nuevas conexiones del mismo IP o
        \emph{SessionID}.
\end{itemize}

%=====================================================================
\section{OAuth 2.0 (Open Authorization)}
%=====================================================================

\textbf{OAuth} es un estándar de \textbf{autorización} ampliamente usado para que las
aplicaciones accedan a recursos de una plataforma \emph{en nombre del usuario}
\textbf{sin que este comparta sus credenciales} con la aplicación. Está detrás de ``Iniciar
sesión con Google/Meta''. Funciona emitiendo \textbf{tokens de acceso temporales y
específicos}; \textbf{OAuth 2.0} es la versión moderna.

\begin{anecbox}{El problema que resuelve}
La forma ``fea'' sería darle a la app mi usuario y contraseña del servicio: eso le daría poder
para borrar la cuenta, cambiar la clave, etc., y mis credenciales quedarían en una base que no
controlo. OAuth nació (Twitter) para que un usuario pueda \emph{delegar} programáticamente a
una aplicación la capacidad de ejecutar cosas en su nombre, \emph{sin entregar credenciales}.
\end{anecbox}

\subsection{El flujo de autorización}
Intervienen el \emph{Resource Owner} (usuario), la \emph{Application}, el
\emph{Authorization Server} y el \emph{Resource Server}.

% ---- Diagrama de secuencia OAuth (recreación slide OAuth 2.0) ----
\begin{figure}[h]
\centering
\begin{tikzpicture}[font=\scriptsize,>=Stealth]
  \def\xa{0} \def\xb{3.3} \def\xc{7.0} \def\xd{10.3}
  \foreach \x/\lab in {\xa/{Resource\\Owner}, \xb/{Application}, \xc/{Authorization\\Server}, \xd/{Resource\\Server}}{
    \node[font=\scriptsize\bfseries,align=center] at (\x,0.4) {\lab};
    \draw[black!40,dashed] (\x,0) -- (\x,-6.2);
  }
  \draw[->] (\xa,-0.4) -- node[above]{1. Access Application} (\xb,-0.4);
  \draw[<-] (\xa,-1.1) -- node[above]{2. Authn \& Request Auth.} (\xb,-1.1);
  \draw[->] (\xa,-1.8) -- node[above]{3. Authn \& Grant Auth.} (\xc,-1.8);
  \draw[<-] (\xb,-2.5) -- node[above]{4. Send Authorization Code} (\xc,-2.5);
  \draw[->] (\xb,-3.2) -- node[above]{5. Code Exchange for Token} (\xc,-3.2);
  \draw[<-] (\xb,-3.9) -- node[above]{6. Issue Access Token} (\xc,-3.9);
  \draw[->] (\xb,-4.6) -- node[above]{7. Request Resource w/ Token} (\xd,-4.6);
  \draw[<-] (\xb,-5.3) -- node[above]{8. Return Resource} (\xd,-5.3);
\end{tikzpicture}
\caption{Flujo de \emph{authorization code}: el usuario autoriza, la app obtiene un
\emph{code}, lo intercambia por un \emph{access token} y con él pide el recurso.
\textit{(Recreación de la diapositiva ``OAuth 2.0''.)}}
\end{figure}

\noindent Puntos clave del flujo:
\begin{itemize}[leftmargin=1.4em,itemsep=1pt]
  \item El usuario nunca comparte con la app las credenciales del servicio; \textbf{delega} el
        control de acceso al \emph{dueño del recurso}, que decide si responde.
  \item El \emph{access token} viaja en un \emph{header} (típicamente \texttt{Bearer}); el
        servidor lo valida en cada requerimiento. El usuario puede \textbf{revocar} el acceso.
\end{itemize}

\subsection{Refresh Token}
Los \emph{access tokens} \textbf{no} duran para siempre (de una hora a unos días). Aparecen
\textbf{dos tokens}: \emph{access} y \emph{refresh}.
\begin{itemize}[leftmargin=1.4em,itemsep=1pt]
  \item Cuando el recurso responde ``token inválido'', el cliente envía el \emph{refresh token}
        y obtiene —\emph{sin} re-preguntar al usuario— un \emph{access} y un \emph{refresh}
        nuevos.
  \item El \emph{refresh token} se usa \textbf{una sola vez y muere}. Esto permite
        \textbf{recuperarse ante un robo}: si un atacante roba ambos tokens, en cuanto uno de
        los dos (atacante o app legítima) usa el refresh, \emph{el otro pierde acceso}.
  \item Así se sostiene la ``sesión infinita'' (Gmail, Mercado Libre\dots) sin el riesgo de un
        token eterno que, si se roba, sería \emph{game over}.
\end{itemize}

\subsection{Scope (alcance) — un mecanismo DAC}
El \textbf{scope} son los permisos específicos que una app solicita para acceder a los datos
del usuario; definen el \textbf{nivel de acceso} que tendrá el token.
\begin{itemize}[leftmargin=1.4em,itemsep=1pt]
  \item Al pedir autorización, la app declara qué \emph{scopes} quiere; eso genera la pantalla
        de consentimiento (``\emph{Get Together} quiere acceder a tu email y a tu calendario'').
  \item El token emitido se \textbf{limita a los scopes otorgados}; cada servicio que lo recibe
        verifica que sea válido y que contenga su propio scope.
\end{itemize}

\begin{anecbox}{OAuth es dual}
La clase siguiente (Autenticación) retoma OAuth: hoy es un protocolo \textbf{dual} que cumple
dos funciones de seguridad. Aquí vimos la primera (\emph{autorización}); la otra
(\emph{autenticación}, vía OpenID Connect) requiere un poco más de teoría.
\end{anecbox}

%=====================================================================
\section*{Lectura recomendada}
\addcontentsline{toc}{section}{Lectura recomendada}
%=====================================================================
\begin{center}
\begin{tcolorbox}[enhanced,colback=itbaBlue!20,colframe=itbaCyanDark,boxrule=1pt,
  arc=3pt,width=0.85\textwidth,halign=center]
\large\textbf{Capítulo 4 \;·\; Capítulo 5 (5.1--5.4) \;·\; Capítulo 6 (6.1--6.2)\\
Capítulo 7 (7.1) \;·\; Capítulo 8 (8.1)}\\[4pt]
\emph{Computer Security: Art and Science}\\[2pt]
Matt Bishop
\end{tcolorbox}
\end{center}
\noindent Los temas de políticas y control de acceso corresponden a los capítulos 4 al 8 del
Bishop, con más ejemplos. Mención especial: el modelo de \textbf{capabilities} (un mecanismo
dual de las ACL, muy útil en sistemas distribuidos a escala gigantesca). Los papers originales
de los modelos: \texttt{Bell76.pdf}, \texttt{Biba75.pdf} y
\texttt{A\_Comparison\_of\_commercial\_and\_military.pdf}.

%---------------------------------------------------------------------
\vfill
\begin{center}\small\itshape\color{black!60}
Apunte integral de la Clase 6 — Políticas de seguridad y Control de acceso · Criptografía y
Seguridad, ITBA.\\
Síntesis de teoría (transcripción + diapositivas) y práctica (diapositivas).
\end{center}

\end{document}
